\documentclass[11pt,a4paper]{article}
\usepackage[margin=24mm]{geometry}
\usepackage[T1]{fontenc}
\usepackage{lmodern,microtype,amsmath,amssymb,amsthm,mathtools,enumitem,xcolor}
\usepackage[colorlinks=true,linkcolor=blue!45!black,urlcolor=blue!45!black]{hyperref}
\newtheorem{theorem}{Theorem}[section]
\newtheorem{lemma}[theorem]{Lemma}
\newtheorem{proposition}[theorem]{Proposition}
\newtheorem{corollary}[theorem]{Corollary}
\theoremstyle{remark}\newtheorem{remark}[theorem]{Remark}
\newcommand{\R}{\mathbb R}\newcommand{\C}{\mathbb C}
\newcommand{\ii}{\mathrm i}\newcommand{\dd}{\,\mathrm d}
\newcommand{\norm}[1]{\left\|#1\right\|}
\newcommand{\ip}[2]{\left\langle #1,#2\right\rangle}
\newcommand{\Rea}{\operatorname{Re}}\newcommand{\Ima}{\operatorname{Im}}
\setlength{\parskip}{3pt}
\title{A mass gap and conserved-momentum geometry\\for negative-energy biharmonic traveling waves}
\author{Research continuation: analytic proofs and independent internal checks}
\date{7 October 2026}
\begin{document}
\maketitle
\begin{abstract}
For the focusing mass-critical biharmonic nonlinear Schr\"odinger equation, we prove a positive gap between the sharp Gagliardo--Nirenberg threshold mass and the mass of any negative-energy traveling profile, in every spatial dimension. The proof uses compactness of optimizing sequences and an elementary rigidity argument for the phase Hessian. It requires neither uniqueness nor real-part nondegeneracy of ground states, and no numerical profile certificate. We also give an exact positive factorization of energy along the entire dilation curve preserving both mass and momentum of a negative-energy wave. These results identify two structural restrictions on coherent escape from collapse. They do not classify all globally bounded solutions.
\end{abstract}
\noindent\textbf{Status.} This is an internally checked research manuscript. The proofs below are self-contained apart from standard Sobolev, multiplier and compactness facts, whose applications are specified. No publication-priority claim or external peer-review claim is made. The main theorems are independent of the computational certificates and multiwave constructions in the earlier project.

\section{Equation and main mass-gap statement}
Fix $N\ge1$ and write
\[
 \alpha=\frac8N,\qquad p=2+\alpha,\qquad r=\frac4N,
 \qquad p=2(r+1).
\]
The equation and profile equation are
\begin{align}
 \ii u_t&=\Delta^2u-|u|^\alpha u,\label{eq:pde}\\
 \Delta^2U+\omega U+\ii v\cdot\nabla U-|U|^\alpha U&=0,
 \qquad u(t,x)=e^{\ii\omega t}U(x-vt).\label{eq:profile}
\end{align}
Here $\omega\in\R$, $v\in\R^N$, and a profile is a nonzero weak solution $U\in H^2(\R^N;\C)$. Set
\[
 M(f)=\norm f_2^2,\quad A(f)=\norm{\Delta f}_2^2,\quad
 E(f)=\frac12 A(f)-\frac1p\norm f_p^p,\quad
 P(f)=\Ima\int\bar f\nabla f.
\]
Let $M_c>0$ be defined by the sharp inequality
\begin{equation}\label{eq:gn}
 \norm f_p^p\le\frac p2\left(\frac{M(f)}{M_c}\right)^r A(f).
\end{equation}
Equivalently, $M_c$ is the mass of a ground state with nonlinear coefficient one. Optimizers can be normalized to have $M=M_c$ and $A=1$.

\begin{theorem}[A uniform gap for all traveling profiles]\label{thm:gap}
For every $N\ge1$ there is $\varepsilon_N>0$ such that every profile of~\eqref{eq:profile} with $E(U)<0$ satisfies
\[
 \boxed{M(U)\ge M_c+\varepsilon_N.}
\]
The gap is uniform in frequency, velocity, translations and phase. No radial or parity assumption is imposed, and no finite-variance hypothesis is needed. The proof is qualitative: it does not supply a numerical value of $\varepsilon_N$.
\end{theorem}

The theorem is an exclusion result; it does not assert existence of negative-energy waves in every dimension. It is stronger than the immediate consequence $M(U)>M_c$ of~\eqref{eq:gn}, since a sequence of strictly larger masses could otherwise converge to $M_c$.

\section{Profile identities at energy regularity}
Every weak profile in~\eqref{eq:profile} belongs to $H^4$. Here are sufficient details to handle small, noninteger $\alpha$ in large dimension. For $N\le8$, Sobolev embedding gives $U\in L^{2(1+\alpha)}$ directly, so $|U|^\alpha U\in L^2$ and~\eqref{eq:profile} gives $H^4$ regularity. For $N>8$, choose $\mu$ so large that
\[
 |\xi|^4-v\cdot\xi+\mu\ge c(1+|\xi|^4).
\]
The inverse $R_\mu=(\Delta^2+\ii v\cdot\nabla+\mu)^{-1}$ maps $L^s$ to $W^{4,s}$ for $1<s<\infty$, by the multiplier theorem. Write
\[
 U=R_\mu(|U|^\alpha U)+(\mu-\omega)R_\mu U.
\]
The second term is in $H^4$, hence in $L^{q_{\rm lin}}$, $q_{\rm lin}=2N/(N-8)$. Starting from $q_0=2N/(N-4)$, the first term improves reciprocal integrability according to
\[
 a_{j+1}=(1+8/N)a_j-4/N,\qquad
 a_0=\tfrac12-\tfrac2N,
\]
until $q=2(1+\alpha)$ is reached. The ideal iterates are $a_j=\tfrac12-(2/N)(1+8/N)^j$, so this takes finitely many steps. Intermediate exponents may be slightly weakened to avoid endpoints. All input exponents $q_j/(1+\alpha)$ exceed one, and $2(1+\alpha)<q_{\rm lin}$. Thus the linear term does not obstruct the iteration. Now the nonlinearity is in $L^2$, and the original equation implies $U\in H^4$.

\begin{lemma}[Pohozaev and current identities]\label{lem:identities}
Every such profile satisfies, with the unitary Fourier transform,
\begin{align}
 v\cdot P(U)&=8E(U),\label{eq:poho}\\
 W(U):=\int |\xi|^2\xi|\widehat U(\xi)|^2\dd\xi&=\frac{M(U)}4v,\label{eq:current}\\
 \omega M(U)&=rA(U)+(8-p)E(U).\label{eq:omega}
\end{align}
For $U=a+\ii b$, with real $a,b$, put
\[
 \mathcal A_U=\Delta^2+\omega-|U|^\alpha.
\]
Then, in quadratic-form notation,
\begin{equation}\label{eq:imaginary}
 \boxed{E(U)=\tfrac14\ip b{\mathcal A_Ub}.}
\end{equation}
\end{lemma}
\begin{proof}
The profile is a critical point of $S=E+\omega M/2-v\cdot P/2$. Mass-preserving dilation scales $E$ and $P$ by $\lambda^4$ and $\lambda$, respectively, giving~\eqref{eq:poho}. Linear phase modulation $e^{\ii k\cdot x}U$ has energy derivative $2W$ and momentum derivative $M$ at $k=0$, giving~\eqref{eq:current}.

These variations do not require weighted assumptions. To justify them at $H^4$, test the weak equation with the cutoff dilation
\[
 (x\chi(x/R))\cdot\nabla U+
 \tfrac12\operatorname{div}(x\chi(x/R))U
\]
and the cutoff phase $\ii x_j\chi(x/R)U$, respectively. Here $\chi$ is smooth, compactly supported and one near the origin. The tests belong to $H^2$. Integration by parts gives the displayed identities in the limit $R\to\infty$: coefficients of the leading terms converge to their constant values, and the additional derivatives of the cutoff are $O(R^{-1})$ or smaller. Their terms are bounded by tails of $|U|^2$, $|\nabla U|^2$, $|\Delta U|^2$, and $|U|^p$, or products controlled by these norms. All tails vanish. The nonlinear integration by parts is valid since $|U|^\alpha U\in L^2$ and $\nabla U\in L^2$.

Taking the real pairing with $U$ gives $A+\omega M-v\cdot P-\norm U_p^p=0$; combine this with $\norm U_p^p=(p/2)A-pE$ to obtain~\eqref{eq:omega}. Finally, the imaginary profile equation is $\mathcal A_Ub=-v\cdot\nabla a$. Since $P(U)=-2\int b\nabla a$, it gives $v\cdot P=2\ip b{\mathcal A_Ub}$, proving~\eqref{eq:imaginary}.
\end{proof}

\section{Phase rigidity without uniqueness or radial symmetry}
The real-valuedness argument below is a standard strict-convexity mechanism; compare the phase-rotation argument in Proposition 3.8 of~\cite{FJMM}. Fourier rearrangement provides further radial symmetry for even powers~\cite{LS}, but radial symmetry is not needed here.

\begin{lemma}[All optimizers have constant phase]\label{lem:phase}
Every complex optimizer of~\eqref{eq:gn} is a constant phase times a real function.
\end{lemma}
\begin{proof}
Let $f=a+\ii b$. With $C=(p/2)M_c^{-r}$, Minkowski,~\eqref{eq:gn} on the real components, and two-term H\"older give
\begin{align*}
 \norm f_p^2
 &\le\norm a_p^2+\norm b_p^2\\
 &\le C^{2/p}\bigl[A(a)^{2/p}M(a)^{1-2/p}
                 +A(b)^{2/p}M(b)^{1-2/p}\bigr]\\
 &\le C^{2/p}A(f)^{2/p}M(f)^{1-2/p}.
\end{align*}
For an optimizer every inequality is an equality. Strict convexity of $L^{p/2}$ gives $b^2=c a^2$ almost everywhere, unless one component vanishes. Apply the same equality argument to $e^{\ii\pi/4}f$. If both signs of $b/a=\pm\sqrt c$ occurred on sets of positive measure, the new two squared components could not be proportional. For $c=1$ they instead have disjoint nonzero supports, which also gives strict Minkowski inequality. Hence one sign holds almost everywhere, proving the claim. The same argument covers any initial constant phase rotation.
\end{proof}

Let $Q$ be a real optimizer normalized by $M(Q)=M_c$ and $A(Q)=1$. Its Euler--Lagrange equation is
\[
 \Delta^2Q+\omega_QQ-|Q|^{2r}Q=0,\qquad \omega_Q=r/M_c>0.
\]

\begin{lemma}[Strict positivity after removing phase]\label{lem:coercivity}
For every such optimizer,
\[
 L_-:=\Delta^2+\omega_Q-|Q|^{2r}\ge0,\qquad
 \ker L_-=\operatorname{span}\{Q\}.
\]
There is $c_Q>0$ such that for every real $b\in H^2$ with $\ip bQ=0$,
\[
 \ip b{L_-b}\ge c_Q\norm b_{H^2}^2.
\]
\end{lemma}
\begin{proof}
We retain $A=A(Q)$ and $M=M(Q)$ in the algebra; $\omega_Q=rA/M$. For real $b$, set $B=M(b)$, $C_1=A(b)$ and
\[
 F(s)=\int(Q^2+s b^2)^{r+1},\qquad s\ge0.
\]
The sharp inequality and equality at $Q$ imply
\begin{equation}\label{eq:comparison}
 F(s)\le (r+1)(1+sB/M)^r(A+sC_1).
\end{equation}
Comparison of right derivatives at zero gives $L_-\ge0$. If $b$ belongs to its form kernel, these first derivatives agree and
\[
 \int |Q|^{2r}b^2=C_1+rAB/M.
\]
Compare the second-order remainders in~\eqref{eq:comparison}. Convexity gives a nonnegative pointwise remainder. Fatou's lemma, or the integral Taylor formula, yields
\begin{equation}\label{eq:Dupper}
 D:=\int |Q|^{2r-2}b^4
 \le (r-1)AB^2/M^2+2BC_1/M.
\end{equation}
Here the integrand at $Q=0$ is interpreted by the limiting second-order remainder: it is infinite if $r<1$ and $b\ne0$, equals $b^4$ if $r=1$, and is zero if $r>1$. Thus when $r<1$ the finite upper bound already forces $b=0$ almost everywhere on that zero set.

Cauchy--Schwarz gives
\[
 D\ge\frac{(C_1+rAB/M)^2}{(r+1)A}.
\]
Combining this with~\eqref{eq:Dupper}, and multiplying by $(r+1)A$, gives
\[
 (C_1-AB/M)^2\le0.
\]
Thus all inequalities are equalities. The Cauchy equality gives $b^2=(B/M)Q^2$ on $\{Q\ne0\}$. Its integral already equals $B$, so $b=0$ on the remaining zero set as well, including the case $r>1$. Moreover $C_1=AB/M$.

The form kernel is a linear space containing $Q$. Applying the same conclusion to $b+Q$ yields $(b+Q)^2=c'Q^2$. Subtracting $b^2=(B/M)Q^2$ gives $bQ=c''Q^2$; since $b$ already vanishes where $Q$ does, $b=c''Q$ almost everywhere. This proves the kernel statement without a maximum principle or a Fourier-positivity argument.

Finally, multiplication by $|Q|^{2r}$ is compact from $H^2$ to $H^{-2}$. Indeed $|Q|^{2r}\in L^{p/(2r)}$, the exponent $p$ is strictly below the $H^2$ critical embedding exponent, local Rellich compactness applies, and the tail is small by H\"older. Consequently $L_-$ is a compact form perturbation of the positive operator $\Delta^2+\omega_Q$. Nonnegativity and the one-dimensional kernel imply the claimed $H^2$ coercivity on the $L^2$ orthogonal complement.
\end{proof}

\section{Compactness and proof of the uniform mass gap}
\begin{lemma}[Normalized near-optimizers are compact]\label{lem:compact}
Suppose $M(f_n)\to M_c$, $A(f_n)=1$, and $E(f_n)<0$. After a subsequence and translations, $f_n$ converges strongly in $H^2$ to a normalized optimizer.
\end{lemma}
\begin{proof}
The sharp inequality sandwiches the nonlinear norm:
\[
 \frac p2<\norm{f_n}_p^p
 \le\frac p2\left(\frac{M(f_n)}{M_c}\right)^r\longrightarrow\frac p2.
\]
In particular $E(f_n)\to0$. Nonvanishing for bounded $H^2$ sequences follows from the local Sobolev inequality and a covering by unit cubes: if the mass on every translated fixed cube tended to zero, the $L^p$ norm would tend to zero. Thus translations yield a nonzero weak limit $Q$, also with local strong convergence and almost-everywhere convergence.

Let $m_1=M(Q)>0$, $a_1=A(Q)>0$, $m_2=M_c-m_1$ and $a_2=1-a_1$. Hilbert norm splitting and the Brezis--Lieb lemma, followed by~\eqref{eq:gn}, give
\[
 M_c^r\le m_1^r a_1+m_2^r a_2,\qquad a_1+a_2=1.
\]
If $0<m_1<M_c$, both coefficients $m_j^r$ are strictly less than $M_c^r$, a contradiction. Hence $m_1=M_c$, and the inequality then forces $a_1=1$. This proves strong $H^2$ convergence and equality in~\eqref{eq:gn}.
\end{proof}

\begin{proof}[Proof of Theorem~\ref{thm:gap}]
If no uniform gap existed, there would be profiles $U_n$ with negative energy and $M(U_n)\downarrow M_c$. Critical scaling
\[
 U(x)\mapsto\lambda^{N/2}U(\lambda x),\qquad
 (\omega,v)\mapsto(\lambda^4\omega,\lambda^3v)
\]
preserves mass and the sign of energy. Normalize $A(U_n)=1$. Lemmas~\ref{lem:compact} and~\ref{lem:phase} give, after translations and a constant phase, $U_n\to Q$ strongly in $H^2$, with $Q$ real. By~\eqref{eq:omega},
\[
 \omega_n\longrightarrow r/M_c=\omega_Q.
\]
Rotate each $U_n=a_n+\ii b_n$ by a further phase tending to zero so that $\ip{b_n}Q=0$. This is possible since $\ip{U_n}Q\to M_c$.

The forms $\mathcal A_n=\Delta^2+\omega_n-|U_n|^{2r}$ converge in operator norm $H^2\to H^{-2}$ to $L_-$. To see this, $U_n\to Q$ in $L^p$ implies $|U_n|^{2r}\to|Q|^{2r}$ in $L^{p/(2r)}$, and H\"older with two $L^p$ factors bounds the form difference. Lemma~\ref{lem:coercivity} therefore gives, for large $n$,
\[
 \ip{b_n}{\mathcal A_n b_n}\ge\tfrac12c_Q\norm{b_n}_{H^2}^2\ge0.
\]
But~\eqref{eq:imaginary} says $E(U_n)=\ip{b_n}{\mathcal A_n b_n}/4<0$. This contradiction proves the uniform gap. No uniform coercivity constant over all ground states is needed: compactness selects one limiting optimizer.
\end{proof}

\begin{corollary}[Arbitrary elliptic quartic anisotropy]
Let $a(\xi)$ be a real homogeneous polynomial of degree four satisfying
\[
 0<c|\xi|^4\le a(\xi)\le C|\xi|^4\qquad(\xi\ne0).
\]
Replace $\Delta^2$ by $a(D)$ and $A(f)$ by $A_a(f)=\int a(\xi)|\widehat f(\xi)|^2\dd\xi$. Define $E_a$ and the sharp threshold mass $M_{c,a}$ accordingly. There is $\varepsilon_{N,a}>0$ such that every $H^2$ profile solving
\[
 a(D)U+\omega U+\ii v\cdot\nabla U-|U|^{8/N}U=0,
 \qquad E_a(U)<0,
\]
satisfies $M(U)\ge M_{c,a}+\varepsilon_{N,a}$.
\end{corollary}
\begin{proof}
The symbol is even, so the kinetic form is real and splits additively over real and imaginary components. Its norm is equivalent to the homogeneous $H^2$ norm, and its homogeneity gives the same dilation identity $v\cdot P=8E_a$. The elliptic regularity argument applies to $a(D)+\ii v\cdot\nabla+\mu$ for large $\mu$. All GN, phase-kernel and compactness arguments above use only these properties. The identity $E_a(U)=\ip b{(a(D)+\omega-|U|^{8/N})b}/4$ is unchanged. The same contradiction proves the result. No smallness of the anisotropy is assumed; the gap is not asserted uniform over all symbols.
\end{proof}

\section{Exact geometry at fixed mass and momentum}
This section applies to every negative-energy profile in every dimension, independently of the mass-gap proof. For such a profile put
\[
 a=\frac{P(U)}{M(U)},\quad m=M(U)|a|^4,\quad
 B=\int\bigl(|a|^2|\xi|^2+2(a\cdot\xi)^2\bigr)|\widehat U(\xi)|^2\dd\xi.
\]
The letters $a,m,B$ here are unrelated to the real/imaginary decomposition in earlier sections. Jensen gives $B\ge3m$.

\begin{theorem}[A strict minimum on the conserved-momentum dilation curve]
Define
\[
 U_\lambda(x)=e^{\ii(1-\lambda)a\cdot x}\lambda^{N/2}U(\lambda x),
 \qquad \lambda>0.
\]
Then $M(U_\lambda)=M(U)$ and $P(U_\lambda)=P(U)$. With $E=E(U)<0$,
\begin{align}\label{eq:factor}
 E(U_\lambda)-E
 =(\lambda-1)^2\bigl[&(B-\tfrac32m-3E)\lambda^2\\
 &+(m-2E)\lambda+(\tfrac12m-E)\bigr].\nonumber
\end{align}
All three coefficients in brackets are positive. In particular,
\[
 E(U_\lambda)-E(U)\ge\bigl(|E(U)|+m/2\bigr)(\lambda-1)^2.
\]
Thus $U$ is the unique energy minimum on this entire curve.
\end{theorem}
\begin{proof}
For $U_{\lambda,k}=e^{\ii k\cdot x}\lambda^{N/2}U(\lambda x)$, direct Fourier expansion gives
\begin{align*}
 E(U_{\lambda,k})={}&\lambda^4E+2\lambda^3 k\cdot W
 +\lambda^2\int\bigl(|k|^2|\xi|^2+2(k\cdot\xi)^2\bigr)|\widehat U|^2\\
 &+2\lambda|k|^2 k\cdot P+\tfrac12M|k|^4.
\end{align*}
The nonlinear energy scales as $\lambda^4$. Equations~\eqref{eq:poho}--\eqref{eq:current} give $a\cdot W=2E$. Substitute $k=(1-\lambda)a$ and factor to get~\eqref{eq:factor}. Positivity follows from $B\ge3m$ and $E<0$.
\end{proof}

Two useful consequences are
\[
 \left.\frac{\dd^2}{\dd\lambda^2}E(U_\lambda)\right|_{\lambda=1}
 =-12E+2B>0,
\]
and
\begin{equation}\label{eq:boost}
 E(e^{-\ii a\cdot x}U)=B-\tfrac32m-3E
 \ge\tfrac32m+3|E|>0.
\end{equation}
So neutralizing a wave's momentum by a linear phase necessarily destroys its negative energy. Ordinary mass-preserving dilation has energy $\lambda^4E$; it does not preserve momentum. A minimum on the particular curve above is not a minimum on the full constraint surface and does not prove orbital stability.

\section{The next dynamical threshold, and a radiation constraint}
Define the distinct dynamical quantity
\[
 m_{\rm ng}=\inf\left\{M(u_0):\begin{array}{l}
 E(u_0)<0,\ P(u_0)=0,\ u\text{ is forward global},\\
 \sup_{t\ge0}\norm{u(t)}_{H^2}<\infty
 \end{array}\right\},
\]
with the infimum of the empty set equal to $+\infty$. The sharp inequality gives only $m_{\rm ng}\ge M_c$. Theorem~\ref{thm:gap} concerns traveling profiles; it does not yet prove $m_{\rm ng}>M_c$. Proving that stronger statement would require control of arbitrary bounded dynamics, including radiation and multiple separating cores.

For comparison, let $m_{\rm tw}$ be the infimum of the masses of all negative-energy traveling profiles. The theorem proves $m_{\rm tw}\ge M_c+\varepsilon_N$, with no assertion that the infimum is attained. If one already has an exact opposite-velocity two-wave solution with profile $U_v$, conservation and decoupling give $m_{\rm ng}\le2M(U_v)$. This upper bound depends on that separate nonlinear construction.

There is a useful universal constraint on a prospective wave-plus-radiation mechanism. Suppose an exact solution were asymptotic to a negative-energy traveling wave plus free radiation $e^{-\ii t\Delta^2}f$, with vanishing cross terms and zero total momentum. For $m_f=M(f)>0$, Jensen gives
\[
 E_{\rm rad}=\tfrac12\norm{\Delta f}_2^2
 \ge\frac{|P(f)|^4}{2m_f^3}.
\]
The constant is sharp as an infimum, approached by narrow Fourier packets. Since $P(f)=-P(U)$, strict negativity of total energy requires
\begin{equation}\label{eq:budget}
 m_f^3>\frac{|P(U)|^4}{2|E(U)|}.
\end{equation}
This is necessary, not a nonlinear existence theorem. For a branch with $|P(U_v)|\asymp|v|$ and $-E(U_v)\asymp v^2$, the lower bound on radiation mass from this budget is of order $|v|^{2/3}$. Thus two full nonlinear cores are not the only possible momentum-balancing mechanism. Constructing an exact radiative wave operator is a substantive next step.

Three objectives now have precise content:
\begin{enumerate}[leftmargin=*,itemsep=3pt]
 \item Determine whether the traveling-wave gap extends to $m_{\rm ng}>M_c$, and identify the least-mass escape mechanism. This is a research question posed here, not an attributed named conjecture.
 \item Construct an infinite-dimensional center-stable family near a negative-energy excited wave, first with uniform orbital control and then with an asymptotic description of radiation. The second objective additionally requires continuous-spectrum and internal-mode analysis.
 \item Prove a local soliton-resolution theorem in a specified regime. Recent one-dimensional mixed-dispersion work explicitly formulates asymptotic stability, branch transitions and soliton-resolution conjectures~\cite{KPRS}. Those are related targets, not consequences of the present all-dimensional mass-gap theorem, and their dimension and parameter restrictions must be preserved.
\end{enumerate}

\begin{thebibliography}{9}
\bibitem{FJMM} A. J. Fern\'andez, L. Jeanjean, R. Mandel and M. Mari\c{s},
\emph{Non-homogeneous Gagliardo--Nirenberg inequalities in $\R^N$ and application to a biharmonic non-linear Schr\"odinger equation},
Journal of Differential Equations 330 (2022), 1--65; Proposition 3.8.
\url{https://arxiv.org/abs/2010.01448}.
\bibitem{LS} E. Lenzmann and J. Sok,
\emph{A sharp rearrangement principle in Fourier space and symmetry results for PDEs with arbitrary order},
International Mathematics Research Notices (2021), 15040--15081; Theorem 2 and its remarks.
\url{https://arxiv.org/abs/1805.06294}.
\bibitem{KPRS} C. Klein, I. Petrenko, S. Roudenko and N. Stoilov,
\emph{Dynamics of solutions in the 1d bi-harmonic nonlinear Schr\"odinger equation},
arXiv:2510.24961v2 (27 February 2026), Section 1.3, Conjecture I.
\url{https://arxiv.org/abs/2510.24961}.
\end{thebibliography}
\end{document}
